% !PRINTED 628
\weberpage{628}
%<<CONT
to the basis $\xi\eta_i$ is effected, we can represent, following the notation (9), (10) of the preceding section, in this way:
\weq{4}{(\xi\eta_i) = Y(\eta_i),}
and accordingly there results, by transition to the basis $\eta_i'$:
\[
(\xi\eta_i') = LYL^{-1}(\eta_i').
\]

The functions $b_1, b_2, \ldots, b_n$ are therefore completely determined by the function $\xi$.

The function
\weq{5}{(-1)^n b_n =
\begin{vmatrix}
y_{1,1}, & y_{1,2}, & \ldots, & y_{1,n} \\
y_{2,1}, & y_{2,2}, & \ldots, & y_{2,n} \\
\cdot & \cdot & \cdot & \cdot \\
y_{n,1}, & y_{n,2}, & \ldots, & y_{n,n}
\end{vmatrix}}
is called the \gesperrt{norm of the function} $\xi$ and is denoted by $N(\xi)$. Concerning it the following theorems hold:

1. \gesperrt{If $\xi$ is not identically zero, then $N(\xi)$ is different from zero.}

For if the determinant of the system (1) vanishes, then rational functions $y_1, y_2, \ldots, y_n$, not all of which vanish, can be determined so that
\[
\xi(y_1\eta_1 + y_2\eta_2 + \cdots + y_n\eta_n) = 0
\]
holds, and this requires, since $\eta_1, \eta_2, \ldots, \eta_n$ is a basis, that $\xi = 0$.

2. \gesperrt{If $\xi$ is a rational function of $z$, then its norm is the $n$th power of this function.}

For if $\xi$ is rational, the equations (1) reduce to the identities $\xi\eta_i = \xi\eta_i$. Accordingly all terms in the determinant (5) vanish with the exception of the diagonal terms, and these all become equal to $\xi$. We express this theorem by the formula:
\[
N(a) = a^n.
\]

3. \gesperrt{If $\xi, \xi'$ are two functions in $\Omega$, then the theorem holds}
\weq{6}{N(\xi\xi') = N(\xi)N(\xi').}

For if, according to (4):
\[
(\xi\eta_i) = Y(\eta_i), \quad (\xi'\eta_i) = Y'(\eta_i),
\]
then
\[
(\xi\xi'\eta_i) = YY'(\eta_i),
\]
and from this formula (6) follows, because the determinant of a composed linear substitution, here $YY'$, is equal to the product of the determinant of the components.
% [sic] printed "Produkt der Determinante der Komponenten" (singular "Determinante"), the sense requires the plural "Determinanten": the determinant of a composed substitution is the product of the determinants of the components
